% Supporting derivations for the companion manuscript, 21 November 2026.
% This is a section proposal, not a standalone document.
% Assumes amsmath, amssymb, amsthm and theorem/proposition/lemma/remark environments.

\section{Hamiltonian gluing and its derivative bundle}

Let $\rho:Z\to\mathbb{CP}^1$ be a holomorphic submersion with
two-dimensional fibers. Write $\zeta$ and $\mu=\zeta^{-1}$ for coordinates
on two direction discs whose union is $\mathbb{CP}^1$. We work on a
neighborhood of a section $L\subset Z$. After choosing fiber coordinates,
write $\xi=(p,q)$ on the northern chart and $\Xi=(P,Q)$ on the southern
chart. A subscript $p$ or $q$ denotes partial differentiation at fixed
$\zeta$. The operator $d_{\mathrm{vert}}$ also differentiates only in
the fiber variables. The line bundle $\mathcal O(k)$ has degree $k$ on
$\mathbb{CP}^1$.

\subsection{The shear reduction}

\begin{proposition}[Fixed-projection symplectic rigidity]
Suppose the first transition coordinate is $P=p/\zeta$. The condition
\begin{equation}\label{eq:audit-symplectic-scaling}
d_{\mathrm{vert}}P\wedge d_{\mathrm{vert}}Q
=\zeta^{-2}dp\wedge dq
\end{equation}
holds if and only if
\begin{equation}\label{eq:audit-shear-transition}
Q=\frac q\zeta+\mathcal T(p,\zeta)
\end{equation}
for a holomorphic function $\mathcal T$ independent of $q$.
\end{proposition}
\begin{proof}
At fixed $\zeta$, the left side of
\eqref{eq:audit-symplectic-scaling} is
$\zeta^{-1}Q_q\,dp\wedge dq$. Thus $Q_q=\zeta^{-1}$.
Integrating in $q$ gives \eqref{eq:audit-shear-transition}; the converse
follows by substitution.
\end{proof}

The proposition concerns the exact prescribed projection $P=p/\zeta$.
It does not exclude different quotient surfaces, different projections,
or more general symplectic gluing.

Set $s(p,\zeta)=\zeta\mathcal T(p,\zeta)$ and remove the standard
line-bundle scaling. The resulting fiber map is
\begin{equation}\label{eq:audit-normalized-shear}
S_\zeta(p,q)=(p,q+s(p,\zeta)).
\end{equation}
At fixed $\zeta$, define
\[
\{f,g\}_{\mathrm{vert}}=f_pg_q-f_qg_p,
\qquad
X_K=\{K,\cdot\}_{\mathrm{vert}}
=-K_q\partial_p+K_p\partial_q.
\]
A local primitive $K$ with $K_p=s$ and $K_q=0$ generates
$S_\zeta$ as its exact time-one flow. It is determined up to a function
of $\zeta$, which generates the zero vector field. The direction parameter
is not differentiated in this Poisson bracket.

\begin{proposition}[The two abelian reductions]
The Hamiltonians depending only on $(p,\zeta)$ form a Poisson-commutative
family. Their normalized derivative transitions are
\begin{equation}\label{eq:audit-normalized-jacobian}
G=DS_\zeta=I+\beta N,
\qquad
N=\begin{pmatrix}0&0\\1&0\end{pmatrix},
\qquad
\beta=s_p=K_{pp}=\zeta\mathcal T_p,
\end{equation}
where $D$ denotes the fiber Jacobian and $I$ the identity matrix.
They lie in the one-dimensional abelian subgroup
$\{I+tN:t\in\mathbb C\}\subset SL(2,\mathbb C)$.
\end{proposition}
\begin{proof}
If $(K_1)_q=(K_2)_q=0$, their vertical Poisson bracket vanishes.
The flows are translations of $q$ by functions of the fixed coordinate
$p$, so $S_s\circ S_t=S_{s+t}$. Differentiation gives
\eqref{eq:audit-normalized-jacobian}. The identity $N^2=0$ gives
$(I+\beta_1N)(I+\beta_2N)=I+(\beta_1+\beta_2)N$.
\end{proof}

The Hamiltonian family is infinite-dimensional; the matrix Lie algebra
$\mathbb C N$ is one-dimensional. Pointwise commutativity, nilpotency of
the chosen matrix, and linearity of a spacetime field equation are
different statements. In particular, the connection $A=\alpha N$, for
a scalar one-form $\alpha$, can have nonzero curvature
$F_A=d\alpha\,N$ despite its abelian color algebra.

If a scalar transition coefficient splits along a complete twistor line
as $\beta=h_0-h_\infty$, with $h_0,h_\infty$ holomorphic on the two
direction discs, then
\[
I+\beta N=(I+h_\infty N)^{-1}(I+h_0N).
\]
Thus the nilpotent matrix representation reduces factorization to scalar
additive splitting. The wave construction supplies explicit factors;
general matrix transitions do not admit this reduction.

\subsection{What the common surface accomplishes}

Let $B$ be a neighborhood of the zero section of $\operatorname{Tot}
\mathcal O(1)$, with coordinates $(p,\zeta)$ and $(P,\mu)$ and transition
$P=p/\zeta$. The shear twistor space and flat twistor space both project
to $B$ by forgetting their second fiber coordinate. Because the matrix
$G$ in \eqref{eq:audit-normalized-jacobian} is independent of $q$, it
defines a bundle $E_B$ on $B$. Pullback gives bundles on both twistor
spaces. This is a geometric common-base transfer, rather than the
identification of the two nonlinear base manifolds.

On a curved wave line and its flat counterpart, the common coordinate is
$p=z-\zeta u=a+\zeta b$ after $a=z$, $b=-u$. Hence the same scalar
factorization reconstructs the curved and flat connections. Equality
of the full pulled-back connection in the wave case is a further fact:
the source connection depends only on the coordinates $u,z$ recovered
independently of direction. It does not follow from equality of a
transition matrix alone.

\subsection{The tangent action of general symplectic gluing}

For a general fiber symplectomorphism $\Phi_\zeta$, both fiber coordinates
may move. An infinitesimal symplectic deformation is locally Hamiltonian.
A general finite symplectomorphism need not be the time-one flow of a
single autonomous holomorphic Hamiltonian; one may instead use a
Hamiltonian isotopy, or work directly with symplectic transition maps.

For a holomorphic vector field $X=X^i(\xi)\partial_{\xi^i}$, define its
tangent lift to the fiber tangent bundle by
\[
\widehat X
=X^i\partial_{\xi^i}
+(\partial_jX^i)v^j\partial_{v^i},
\]
where $\xi^1=p$, $\xi^2=q$, and $v^1,v^2$ are tangent-vector components.
Repeated indices $i,j$ range over $1,2$ and are summed. Direct
differentiation gives
\begin{equation}\label{eq:audit-tangent-lift-bracket}
[\widehat X,\widehat Y]=\widehat{[X,Y]}.
\end{equation}
In particular,
$[\widehat X_K,\widehat X_L]=\widehat X_{\{K,L\}_{\mathrm{vert}}}$.
However, $X\mapsto DX$ alone is not a homomorphism into pointwise
matrix-valued functions. Its exact identity is
\[
D[X,Y]=X(DY)-Y(DX)+(DY)(DX)-(DX)(DY).
\]
Here $X(DY)$ means applying the derivation $X$ to every matrix entry.
The finite counterpart is the chain rule
\begin{equation}\label{eq:audit-chain-rule}
D(\Phi\circ\Psi)(\xi)=D\Phi(\Psi(\xi))D\Psi(\xi).
\end{equation}
The evaluation-point transport disappears in the shear reduction because
every map preserves $p$ and every Jacobian depends only on $p$.

For example, $K_1=p^3/3$ and $K_2=q^3/3$ give
\[
\{K_1,K_2\}_{\mathrm{vert}}=p^2q^2,
\qquad
[X_{K_1},X_{K_2}]
=-2p^2q\,\partial_p+2pq^2\,\partial_q.
\]
This algebraic example illustrates the additional terms; it is not by
itself a complete twistor-space construction. The direction dependence,
patch domains and normal bundles must still satisfy the geometric
hypotheses below.

\section{Derivative bundles beyond the shear family}

\begin{proposition}[General patchwise tangent-line charts]
Let $\rho:Z\to\mathbb{CP}^1$ be a holomorphic submersion with a complete
local four-parameter family of sections $L_x$, $x\in M$, whose normal
bundles are $\mathcal O(1)^{\oplus2}$. On a direction disc choose fiber
coordinates $w^A$, $A=0,1$, and write the sections as
$w^A=w^A(x,\zeta)$. Then
\begin{equation}\label{eq:audit-general-tangent-chart}
y^{A0'}=w^A-\zeta\partial_\zeta w^A,
\qquad
y^{A1'}=\partial_\zeta w^A,
\qquad \zeta\longmapsto\zeta
\end{equation}
defines a local biholomorphism from the curved correspondence space
to flat correspondence space. It preserves the alpha distributions.
At a fixed source point its flat spacetime image is a possibly stationary
null curve with indicated primed direction $(-\zeta,1)$.
\end{proposition}
\begin{proof}
Completeness identifies $T_xM$ with $H^0(L_x,N_{L_x/Z})$. Differentiating
the section in a parameter direction gives the normal section whose
local components are $\delta w^A$. At a fixed $\zeta$, the differential
of \eqref{eq:audit-general-tangent-chart} is an invertible triangular
linear transformation of its first jet $(\delta w,\partial_\zeta\delta w)$.
The first-jet evaluation
\[
H^0(L_x,N_{L_x/Z})\longrightarrow J^1_\zeta N_{L_x/Z}
\]
is an isomorphism: its kernel is
$H^0(L_x,N_{L_x/Z}(-2\zeta))=H^0(\mathcal O(-1)^{\oplus2})=0$,
and both spaces have dimension four. Here $-2\zeta$ denotes vanishing
to order two at the specified point, and $J^1_\zeta$ is the space of
value and first-derivative jets there. Thus the four spacetime
coordinates in \eqref{eq:audit-general-tangent-chart} have invertible
parameter differential; retaining $\zeta$ gives a correspondence-space
local biholomorphism.

The curved twistor projection has coordinates $(\zeta,w^A(x,\zeta))$.
The flat twistor projection has coordinates
$(\zeta,y^{A0'}+\zeta y^{A1'})$, which are exactly the same functions
under \eqref{eq:audit-general-tangent-chart}. Their kernels, the alpha
distributions, therefore correspond.

Finally, at fixed $x$,
\[
\frac{\partial y^{AA'}}{\partial\zeta}
=\partial_\zeta^2w^A\,(-\zeta,1)^{A'}.
\]
This has rank at most one, hence is null for the flat determinant metric
$\eta=2(dy^{00'}dy^{11'}-dy^{01'}dy^{10'})$, and it may vanish.
\end{proof}

The proposition uses no common surface $B$ and no relative symplectic
form. It therefore separates the local existence of generalized
isothermal charts from the much stronger full-connection descent
requirement. The direction disc is essential; the statement does not
glue the inverse charts over complete target projective fibers.

\begin{proposition}[Line triviality of the twisted vertical tangent]
Let $\rho:Z\to\mathbb{CP}^1$ be a holomorphic submersion, and let
$L\subset Z$ be a section whose normal bundle is
$N_{L/Z}\simeq\mathcal O(1)\oplus\mathcal O(1)$. Then
\[
E=T_{Z/\mathbb{CP}^1}\otimes\rho^*\mathcal O(-1),
\qquad T_{Z/\mathbb{CP}^1}=\ker d\rho,
\]
restricts to the trivial rank-two bundle on $L$. If $Z$ also has a
fiberwise nondegenerate relative two-form with values in
$\rho^*\mathcal O(2)$, then $E$ carries a nondegenerate holomorphic
alternating form with scalar values and has structure group $SL(2,\mathbb C)$.
\end{proposition}
\begin{proof}
Since $\rho|_L$ is an isomorphism, $TL$ is complementary to
$T_{Z/\mathbb{CP}^1}|_L$ inside $TZ|_L$. Projection to the quotient by
$TL$ identifies the latter bundle with $N_{L/Z}$. Twisting by
$\mathcal O(-1)$ therefore gives $E|_L\simeq\mathcal O^{\oplus2}$.
The relative two-form is a section of
$\Lambda^2T_{Z/\mathbb{CP}^1}^*\otimes\rho^*\mathcal O(2)
=\Lambda^2E^*$, proving the last assertion.
\end{proof}

The nonlinear-graviton/hyperk\"ahler correspondence applies when the
twistor space has this relative symplectic structure and a suitable
four-complex-dimensional local family of sections with normal bundle
$\mathcal O(1)^{\oplus2}$. Its real form additionally requires the
appropriate reality structure. This standard correspondence supplies
the broader gravitational setting; the preceding line-triviality proof
itself is elementary and uses no field equation.

\begin{lemma}[Openness of line triviality]
Let $\mathcal V$ be a holomorphic rank-$r$ vector bundle on
$U\times\mathbb{CP}^1$, where $U$ is a connected complex neighborhood of
$0$, and suppose $\mathcal V|_{\{0\}\times\mathbb{CP}^1}$ is trivial.
After shrinking $U$, every fiber restriction is trivial and
\[
\mathcal V\simeq p_U^*(p_{U*}\mathcal V),
\]
where $p_U$ is projection to $U$.
\end{lemma}
\begin{proof}
The degree of the restriction is locally constant and equals zero.
At $0$, the bundle $\mathcal V(-1)$ has no global holomorphic sections.
Upper semicontinuity for the proper projection $p_U$ implies
$H^0(\mathbb{CP}^1,\mathcal V_y(-1))=0$ for nearby $y$.
By the Birkhoff--Grothendieck theorem,
$\mathcal V_y\simeq\bigoplus_{j=1}^r\mathcal O(k_j)$.
The vanishing says every $k_j\leq0$, while degree zero says
$\sum_jk_j=0$. Therefore all $k_j=0$.
The dimensions $h^0=r$, $h^1=0$ are now constant. Direct image and base
change give a rank-$r$ bundle $p_{U*}\mathcal V$, and its evaluation map
is an isomorphism on every fiber, hence an isomorphism of bundles.
\end{proof}

\begin{proposition}[A chart-dependent auxiliary flat Ward bundle]
Assume that a neighborhood of a section $L\subset Z$ admits two
fiber-coordinate charts over the direction discs, with $L$ represented
by $(p,q)=0$ and $(P,Q)=0$. Suppose the overlap has the form
\begin{equation}\label{eq:audit-general-transition}
\begin{pmatrix}P\\Q\end{pmatrix}
=\zeta^{-1}\Phi_\zeta(p,q),
\qquad \Phi_\zeta(0)=0,
\end{equation}
where $\Phi_\zeta$ is holomorphic and locally invertible on a uniform
neighborhood of zero over a slightly smaller closed overlap annulus.
Suppose $N_{L/Z}\simeq\mathcal O(1)^{\oplus2}$. Then the matrix
\[
G(p,q,\zeta)=D\Phi_\zeta(p,q)
\]
can also be used as the transition of a holomorphic rank-two bundle
$E_{\mathrm{flat}}$ on a sufficiently small neighborhood of the zero
section in flat twistor space
$Z_{\mathrm{flat}}=\operatorname{Tot}(\mathcal O(1)\oplus\mathcal O(1))$.
It is trivial on every flat twistor line over a sufficiently small
spacetime neighborhood and therefore reconstructs a local holomorphic
ASD connection. If the transition preserves the relative symplectic
form, the bundle and connection can be taken with structure group
$SL(2,\mathbb C)$.
\end{proposition}
\begin{proof}
In the flat charts the base transition is
$(P,Q)=\zeta^{-1}(p,q)$. Regard $G(p,q,\zeta)$ as a matrix function in
the northern coordinates of this flat overlap, rather than using the
curved base transition \eqref{eq:audit-general-transition}. After
shrinking the fiber domains, it is holomorphic and invertible throughout
the overlap. There are only two patches; gluing by $G$ and $G^{-1}$
therefore defines a holomorphic bundle without a further triple-overlap
condition.

Its restriction to the zero section has transition $G(0,0,\zeta)$.
The same matrix is the transition of
$T_{Z/\mathbb{CP}^1}|_L\otimes\mathcal O(-1)$, so the preceding
proposition makes that restriction trivial.

Write $y=(a,b,c,d)$ for flat spacetime coordinates, and use the flat
incidence map
\[
q_{\mathrm{flat}}(y,\zeta)
=(p,q,\zeta)=(a+\zeta b,c+\zeta d,\zeta).
\]
For sufficiently small $y$, its complete projective fibers lie in the
chosen twistor neighborhood. Apply the openness lemma to
$\mathcal V=q_{\mathrm{flat}}^*E_{\mathrm{flat}}$ on
$U\times\mathbb{CP}^1$. Every fiber is then trivial. The pullback from
twistor space has a canonical flat partial connection along the kernel
of $dq_{\mathrm{flat}}$, the flat alpha distribution. The standard
Ward reconstruction therefore gives an ordinary ASD connection on $U$.
Equivalently, in a frame from $p_{U*}\mathcal V$, its alpha-direction
coefficients are sections of $\mathcal O(1)$ on each projective fiber,
so they are linear in the homogeneous direction spinor and are the
restrictions of a unique spacetime connection.

If the relative symplectic condition holds, then
$\det D\Phi_\zeta=1$. Hence the determinant bundle has its canonical
trivialization, and reconstruction respects it after choosing a
determinant-one spacetime frame.
\end{proof}

\begin{remark}[Scope of the auxiliary construction]
The proposition is an existence statement relative to the specified
two-chart presentation. It does not assert invariance under changes of
the curved fiber coordinates. Indeed, the southern derivative of a
coordinate change is evaluated at the curved southern point
$\zeta^{-1}\Phi_\zeta(p,q)$, while on flat twistor space that point is
$\zeta^{-1}(p,q)$. The resulting factors need not extend as an ordinary
flat-twistor bundle gauge transformation. The construction also does not
identify its reconstructed connection with the full source spin
connection pulled back through a generalized chart.

The assertion that some auxiliary flat ASD bundle exists is therefore
strictly weaker than geometric full-connection descent. The shear model
has the additional common-surface mechanism and an explicit equality of
full connection one-forms. Those properties require separate proofs in
any extension.
\end{remark}

% Bibliographic entries for integration into the main manuscript.
% Known theorems are cited below; the elementary propositions above are
% derivations from their stated hypotheses, not novelty claims.
\begin{thebibliography}{9}
\bibitem{HKLR}
N.~J. Hitchin, A. Karlhede, U. Lindstr\"om and M. Ro\v cek,
\emph{Hyperk\"ahler metrics and supersymmetry},
Commun. Math. Phys. \textbf{108} (1987), 535--589.
\href{https://doi.org/10.1007/BF01214418}{doi:10.1007/BF01214418}.
Theorem 3.3 states the twistor reconstruction with reality data;
the holomorphic local form omits the reality requirement.

\bibitem{Grothendieck}
A. Grothendieck,
\emph{Sur la classification des fibr\'es holomorphes sur la sph\`ere de Riemann},
Amer. J. Math. \textbf{79} (1957), 121--138.
\href{https://doi.org/10.2307/2372388}{doi:10.2307/2372388}.

\bibitem{Grauert}
H. Grauert,
\emph{Ein Theorem der analytischen Garbentheorie und die Modulr\"aume
komplexer Strukturen}, Publ. Math. IH\'ES \textbf{5} (1960), 5--64;
correction, \textbf{16} (1963), 35--36.
\href{https://doi.org/10.1007/BF02684746}{doi:10.1007/BF02684746}.

\bibitem{Ward}
R.~S. Ward, \emph{On self-dual gauge fields},
Phys. Lett. A \textbf{61} (1977), 81--82.
\href{https://doi.org/10.1016/0375-9601(77)90842-8}{doi:10.1016/0375-9601(77)90842-8}.
\end{thebibliography}
